\documentclass[aps,prl,reprint,superscriptaddress,nofootinbib,floatfix]{revtex4-2}
\usepackage[utf8]{inputenc}
\usepackage[T1]{fontenc} 

\usepackage{amsmath}
\usepackage{amssymb}
\allowdisplaybreaks

\usepackage{graphicx}
\usepackage{dcolumn}
\usepackage{xcolor}
\usepackage{float}
\usepackage{bm}
\usepackage{comment}
\usepackage{tabularx}
\usepackage{longtable}
\usepackage{booktabs}

\usepackage[italicdif
f]{physics}
\usepackage[colorlinks=True,allcolors=blue]{hyperref}



%%%%%%%%%%%%%%%%%%%%%%%%%%%%%%%%%%%%
\def\DeptOfPhyIITR{Department of Physics,
Indian Institute of Technology Roorkee, Roorkee 247667, India} 

\def\SchOfMathPhyXMUM{School of Mathematics and Physics,
Xiamen University Malaysia, 
43900 Sepang, Malaysia}

\def\InstOfThPhyAstroUG{Institute of Theoretical Physics and Astrophysics,
%Faculty of Mathematics, Physics and Informatics,
University of Gdańsk,
80-308 Gda\'{n}sk, Poland}

\def\CentOfPhotQuCommTechIITR{Centre for Photonics and Quantum Communication Technology, 
Indian Institute of Technology Roorkee, Roorkee 247667, India}
%%%%%%%%%%%%%%%%%%%%%%%%%%%%%%%%%%%%



\begin{document}


%%%%%%%%%%%%%%%%%%%%%%%%%%%%%%%%%%%%
%%%%%%%%%%%%%%%%%%%%%%%%%%%%%%%%%%%%
\title{Testing Modified Gravity with Coupled Microspheres}


%%%%%%%%%%%%%%%%%%%%%%%%%%%%%%%%%%%%
\author{Ankit Kumar}
\email{kumar.ankit.vyas@gmail.com}
\affiliation{\DeptOfPhyIITR}

%%%%%%%%%%%%%%%%%%%%%%%%%%%%%%%%%%%%
\author{Yen-Kheng Lim}
\affiliation{\SchOfMathPhyXMUM}

%%%%%%%%%%%%%%%%%%%%%%%%%%%%%%%%%%%%
\author{P. Arumugam}
\affiliation{\DeptOfPhyIITR}
\affiliation{\CentOfPhotQuCommTechIITR}

%%%%%%%%%%%%%%%%%%%%%%%%%%%%%%%%%%%%
\author{Tom Z\l{}o\'{s}nik}
\affiliation{\InstOfThPhyAstroUG}

%%%%%%%%%%%%%%%%%%%%%%%%%%%%%%%%%%%%
\author{Tomasz Paterek}
\affiliation{\SchOfMathPhyXMUM}
\affiliation{\InstOfThPhyAstroUG}


%%%%%%%%%%%%%%%%%%%%%%%%%%%%%%%%%%%%
%%%%%%%%%%%%%%%%%%%%%%%%%%%%%%%%%%%%
\begin{abstract}
While a wide variety of astrophysical and cosmological phenomena suggest the presence of Dark Matter, all evidence remains via its gravitational effect on the known matter. As such, it is conceivable that this evidence could be explained by a modification to gravitation or concepts of inertia. Various formulations of modified gravity exist, each giving rise to several non-canonical outcomes. This motivates us to propose experiments searching for departures from (quantum) Newtonian predictions in a bipartite setting with gravitational accelerations $\lesssim 10^{-10}$ m/s$^2$, i.e., where the effective force needs to be stronger than Newtonian to account for the Dark Matter effects. Since quantum particles naturally source weak gravitation, their non-relativistic dynamics offers opportunities to test this small acceleration regime. We show that two nearby quantum particles accumulate significantly larger entanglement in modified gravity models, such as the Modified Newtonian Dynamics (MOND). We demonstrate how the temperature can be fine-tuned such that these effects are certified simply by witnessing the entanglement generated from uncorrelated thermal states, eliminating the need for precise noise characterization. Our second regime of interest is bipartite dynamics in the presence of external gravitational fields. The MOND model does not adhere to the Strong Equivalence Principle, leading to anisotropic mutual attraction dependent on the orientation of the symmetry axis with respect to the external field. We calculate the temporal resolution required to detect this phenomenon on Earth.
\end{abstract}



%%%%%%%%%%%%%%%%%%%%%%%%%%%%%%%%%%%%
\maketitle
%%%%%%%%%%%%%%%%%%%%%%%%%%%%%%%%%%%%


The Newtonian limit of General Relativity is very successful on the scale of the solar system.
For example, balancing the centrifugal and the gravitational forces of objects in approximately circular orbits around the Sun implies orbital velocity falls as the square root of the distance. 
This is famously known as Keplerian decline and has been observed to hold for all planets~\cite{NASA-PlanetFactSheet-rel2Earth}.
Spiral galaxies have a lot in common with the solar system. 
Most of their mass is also concentrated towards the center, but the stars do not show any asymptotic Keplerian decline ~\cite{book-DM-Sanders2010}. 
Their orbital speeds generally do not fall, and the rotation curves saturate~\cite{Andromeda-Babock1939}.
Consequently, the stars in the outer regions appear to be orbiting so fast that they cannot be gravitationally bound.
This is not happening, and hence there seems to be more gravity than expected based on the known mass at the center of spiral galaxies. 
This is a prime example of the Dark Matter (DM) effect. The name originates in the proposal of the existence of an invisible matter distributed throughout the galaxies~\cite{ComaCluster-Zwicky1937}, generating an extra gravitational pull that balances the centrifugal force.
Despite being the most widely accepted explanation and with evidence appearing even on the largest cosmological scales, it has not been directly detected or confirmed by any experiment~\cite{DM-Review-Bertone2018}, hence the continued interest in alternative solutions.
A plausible route involves modifications to our present understanding of gravity.
Modified Newtonian Dynamics (MOND) is one such proposal where, without invoking DM, Newton's second law and/or the law of universal gravitation is modified to account for DM effects in galaxies~\cite{MOND-Milogram1983}.
While the experiments we propose are independent of any concrete formulation of alternative models, for quantitative statements we will mostly follow the parameters present in the MOND theory. We therefore begin with more details about it.

A general form of MOND, which encompasses a wide variety of proposed variants of the model, generalizes Newton's second law to 
\begin{equation}
\nu\bigg(\frac{|a|}{a_{0}}\bigg)\vec{a} = -\vec{\nabla}\Phi + \vec{F}/m,
\label{eq:DiffMOND-1}
\end{equation}
where $\vec{F}$ represents the sum of non-gravitational forces on an object of mass $m$, $\vec a$ is its acceleration, and $a_0$ is the acceleration scale where the generalization is to take place. The potential $\Phi$ belongs to the set of potentials $\{ \phi_{1}, \phi_{2}, \dots \}$ or it could be their linear combination.
Collectively, the potentials obey a system of potentially nonlinear and coupled Poisson equations:
\begin{equation}
\sum_{j}\vec{\nabla}\cdot\bigg[\mu_{ij}\bigg(\frac{\vec{\nabla}\phi_{1}}{a_{0}}, \frac{\vec{\nabla}\phi_{2}}{a_{0}}, \dots \bigg)\vec{\nabla} \phi_{j}\bigg] = 4\pi G_{i} \rho,
    \label{eq:DiffMOND-2}
\end{equation}
where $\rho$ is the density of matter, and 
$(\nu,\mu_{(ij)})$ are called the interpolating functions which govern the transition from Newtonian dynamics to the modified regime. 
The theory does not determine their exact functional form, but consistency with experimental data fixes the limits. For example, in the Lagrangian-based AQUAL model~\cite{Lag4MOND-Bekenstein1984} $\nu=1$ and there exists a single potential $\Phi$ with interpolating function $\tilde{\mu}$. In the limiting cases 
$\tilde \mu(x) \to x$ for $x\ll 1$ to explain galactic rotation curves and $ \tilde \mu(x) \to 1$ for $x\gg 1$ to recover Newtonian dynamics in regimes of stronger gravity.
Note that even a single massive object may source multiple potentials $\phi_{(i)}$.
In the Appendix we show how many existing models are special cases of the above expressions.
Notably, although each model is non-relativistic, some exist as limits of relativistic completions. 
A recent example of such a completion, building on antecedent work~\cite{TeVeS-Bekenstein2004}, was shown to be consistent with data from the large-scale cosmic structure and inhomogeneities of the cosmic microwave background~\cite{RMOND-Zlosnik2021}.
Although there does not exist a unique model for the MOND paradigm, we will see that many of them and potentially other modified gravity theories predict deviations from Newtonian gravity that may be observable in bipartite precision experiments.


%%%%%%%%%%%%%%%%%%%%%%%%%%%%%%%%%%%%
A wide array of DM phenomenology in galaxies is associated with acceleration due to gravity below $a_0 \approx 1.2 \times 10^{-10}$ m/s$^2$~\cite{MOND-Milogram1983,MONDinMilkyWay-Famaey2005,MOND-Gentile2011}. Therefore, our first regime of interest is that of small accelerations. 
We will show how the entanglement dynamics between two microspheres senses the force gradient between them, and design tests where a simple act of entanglement witnessing reveals non-canonical interaction. 
The modified gravity models have a stronger force gradient than Newtonian in the discussed limit, thereby leading to a stronger entanglement which is robust against noisy measurements.
Such an experiment could be done using the masses recently cooled in Vienna~\cite{Sciene-Aspelmeyer2020} which, when separated by a distance of a few times their radius, generate an acceleration deep into the discussed relevant regime.
Our calculations assume the particles are in free space devoid of all external gravitational fields.
Although this might seem unrealistic, there may be places in the solar system where the net gravitational acceleration is in the MOND regime~\cite{MOND-Bekenstein2006,MOND-LISA-Magueijo2012}.
Furthermore, it is also instructive to perform this experiment with freely falling microspheres next to Earth's surface. 
An observation of non-canonical effects in agreement with our calculations would indicate that the modified gravity adheres to the Strong Equivalence Principle (SEP), i.e., the internal dynamics stays the same in all uniformly accelerated frames (up to corrections due to tidal effects).
Another possibility motivating the experiment on Earth is the modification of the law of inertia, in which case the non-gravitational forces could be used to balance gravitational accelerations in order to achieve the required regime of small accelerations~\cite{MOND-Ignatiev2007,MOND-Ignatiev2015,MOND-Ignatiev2009,ModGrav-TorsionPend-Timberlake2021}.

%%%%%%%%%%%%%%%%%%%%%%%%%%%%%%%%%%%%
On Earth, various MOND models predict the so-called external field effect (EFE).
The EFE is a consequence of nonlinear equations of motion, putting the MOND theory in conflict with the SEP. 
Therefore, the second regime we study is the internal dynamics of two masses embedded in an external gravity where the EFE constrains the bipartite gravity to Newtonian, albeit with an anisotropic Newton's constant~\cite{EFE-Blanchet2011}.
This anisotropy is quantified by the alignment with the external field and the derivative of the MOND interpolating function, which is a largely unspecified parameter of the model.
We study the motion of two classical masses arranged in parallel and perpendicular to the Earth's gravity and derive the temporal precision required to resolve the deviations from Newtonian dynamics.


%%%%%%%%%%%%%%%%%%%%%%%%%%%%%%%%%%%%
Many proposals have already been put forward to test the predictions of MOND. Torsion pendula experiments show an agreement with Newton's second law down to accelerations of $5 \times 10^{-14}$ m/s$^2$~\cite{NewtonLaw-TorsionPend-Gundlach2007}, and with Newton's law of universal gravity down to $2 \times 10^{-12}$ m/s$^2$~\cite{TestNewton-Little2014}.
It has been argued that the STEP test of the equivalence principle~\cite{EquiPri-Worden1978} can also serve as a tool to verify the compliance of MOND theories with the SEP~\cite{MOND-SEPwSTEP-Pereira2016}.
Note that the accelerations mentioned in these experiments are with respect to the local laboratory, and the net acceleration due to heavenly sources is well above $a_0$. Therefore, it remains unclear to what extent departures from Newtonian dynamics are to be expected.



%SQUID-based measurements of the interaction between magnetically levitated resonators have demonstrated the potential to probe the deviations from Newtonian gravity and admit spatial resolution of a femtometer~\cite{ModGrav-TorsionPend-Timberlake2021}. 
%In this work, we treat MOND as a motivation to consider specific experimental setups in more general context of modified gravity.


%%%%%%%%%%%%%%%%%%%%%%%%%%%%%%%%%%%%%%%%%%%%%%%%%%%%%%%%%%%%%%%%%%%%%%%%
\emph{Small accelerations.}
The driving idea behind modifying standard theories of gravity for small accelerations is the following: just like Newtonian gravity is an approximation of General Relativity when the gravitational field is not too strong, it might also be an approximation of an underlying theory when the Newtonian accelerations are not too small.
In order to match the stellar rotation curves, the modification has to involve accelerations smaller than $a_0 \approx 1.2 \times 10^{-10}$ m/s$^2$, and the simplest solution is to demand that gravitational force in this regime scales inversely to the distance.
Accordingly, in the deep MOND regime the modified gravitational potential at a position $\vec r_2$ due to a mass $m$ located at position $\vec r_1$ is given by $\Phi = \sqrt{Gma_0} \ \ln( \abs{\vec r_2 - \vec r_1})$. 
Note that this violates the law of equal and opposite action and reaction: the force on a mass $m_2$ due to $m_1$ is $\sim m_2\sqrt{m_1}$ whereas the force on mass $m_1$ due to $m_2$ is $\sim m_1\sqrt{m_2}$.
This issue is rectified if one solves the full nonlinear Poisson-like equation governing the bipartite dynamics in MONDian gravity.
The resultant gravitational potential energy of two identical particles of mass $m$ is given by~\cite{MOND2Body-Milgrom2014,MOND2Body-Zhao2010,MOND2Body-Milgrom1986}:
\begin{equation}
V(\vec r_1, \vec r_2) =  \frac{4}{3} \qty( \sqrt{2} - 1 ) \ m \sqrt{G m a_0} \ \ln(L+r) ,
\label{eq:ForceIdenticalMasses}    
\end{equation}
where $r$ is the \emph{relative displacement} of the two masses from their initial separation $L$.
This is very different from the usual Newtonian potential and underlies the differences in observable quantities.
In particular, we focus on quantum entanglement between small objects.



%%%%%%%%%%%%%%%%%%%%%%%%%%%%%%%%%%%%%%%%%%
Entanglement dynamics of two nearby quantum masses in empty space, starting with uncorrelated thermal states, has been studied in detail in  Refs.~\cite{QuantGrav-Tanjung2020,QuantGrav-MesEnt-Qvarfort2020,QuantGrav-Ankit2022}. The system gets entangled because of the position dependence of the gravitational force: the parts of the wave packet closer to each other are attracted more than the parts further away. Accordingly, different momenta are generated at different positions, producing entanglement as time passes.
Notably, the methods of Ref.~\cite{QuantGrav-Ankit2022} apply to any central interactions without the need for an explicit construction for the Hamiltonian.
The displacement-to-separation ratio $r/L$ is always small, and hence we expand Eq.~\eqref{eq:ForceIdenticalMasses} in a binomial series. Note that the cubic term is essential in entanglement dynamics only when there is a significant relative motion between the two particles~\cite{QuantGrav-Ankit2022}, and hence we truncate at the quadratic term only:
\begin{equation}
V  \approx  \frac{4}{3} \qty( \sqrt{2} - 1 ) \ m \frac{\sqrt{Gma_0}}{L} \qty( \ln(L) +  \frac{r}{L} - \frac{r^2}{2L^2} ) .
\end{equation}
The system admits Gaussianity at all times, and we quantify entanglement through the logarithmic negativity of the bipartite covariance matrix~\cite{PPTCrit-CV-Simon2000,Negativity-Gaussian-Adesso2005,Negativity-Vidal2002,Negativity-CV-Serafini2004}.
The covariance matrix starting from the uncorrelated ground states of individual harmonic traps, each with frequency $\omega_0$, is derived in exact closed form in Ref.~\cite{QuantGrav-Ankit2022} and for completeness presented in the Appendix.
If instead one starts with thermal states, characterized by the average phonon number $\bar{n} = 1/ \qty[\exp(\hbar\omega_0/k_BT)-1]$, the covariance matrix is given by $\qty(2\bar{n}+1)$ times the one obtained for the initial ground state.
Furthermore, the covariance matrix is entirely determined by the frequencies $ \omega_0 $ and $ \omega $, where the latter is a function of the interaction between the masses.
For the standard Newtonian gravity one finds $\omega^2 = \omega_N^2 = 4Gm/L^3$~\cite{QuantGrav-Ankit2022}, and for the MONDian gravity we obtain $\omega^2 = \omega_M^2 = 8(\sqrt{2}-1)\sqrt{Gma_0}/3L^2$.
A short algorithm to find $\omega$ for arbitrary central forces is given in the Appendix.


%%%%%%%%%%%%%%%%%%%%%%%%%%%%%%%%%%%%
The choice of microspheres is now obtained as follows. First of all, the requirement of acceleration being smaller than $a_0$, i.e., $Gm/L^2 < a_0$, for a typical material of density $\rho_0 \sim 10$ g/cm$^3$, is satisfied by spheres of radius $R_0 \lesssim 100 \ \mu$m. 
However, not all such configurations lead to higher entanglement in the modified gravity because  entanglement is driven by the force gradient, not the force~\cite{QuantGrav-Ankit2022}. 
The MONDian entanglement is larger than Newtonian when $\omega_M > \omega_N$. 
This condition is more demanding than the previous one and requires spheres of radius $R_0 \lesssim 10 \ \mu$m.
Entanglement dynamics for two Osmium spheres (densest naturally occurring material) with $R_0 = 0.25 \ \mu$m separated by $L = 2.5R_0$
is presented in Fig.~\ref{fig:entMOND}.
This configuration admits an internal acceleration of $\approx 2.5 \times 10^{-13}$ m/s$^2$, readily in the regime of our interest.
The initial (Gaussian) state is prepared by cooling the masses in harmonic traps of frequency $\omega_0 = 25$ kHz.
The left panel shows the ideal case of $T=0$, whereas the right panel is at temperature $T = 0.05 \ \mu$K.
The MONDian entanglement is much higher than what is accumulated by Newtonian gravity.
Note that the shaded region within the two curves corresponds to other possible approaches to modified gravity predicting a force gradient between Newtonian and MONDian.
We present the calculation for finite temperature because a detectable entanglement is accumulated exclusively in the modified gravity theories.
Furthermore, note that the experimenter is not required to quantify entanglement as mere entanglement witnessing within a time window of $1 \lesssim t \lesssim 2$ seconds would indicate the presence of modified gravity.
Clearly, for any configuration, the temperature can be adjusted such that a detectable entanglement is accumulated exclusively in modified gravity models.


\begin{figure}[!t]
\centering
\includegraphics[width=\linewidth]{entMOND_Os_R0.25e-6_L2.5R_w02.5e4-paper.pdf}
\label{fig:entMOND_Os}
\caption{Gravitational entanglement between two identical Osmium spheres of radius $0.25 \ \mu$m initially separated by $2.5$ times their radius. The initial state is a natural Gaussian prepared through cooling in harmonic traps of frequency $25$ kHz.
Entanglement $E$ is measured by logarithmic negativity (note logarithmic vertical scale), $t$ denotes the time, and $T$ is the temperature.
The grey shaded region, where $E < 0.01,$ signifies entanglement not accessible with current technological precision~\cite{Science-Palomaki2013}.}
\label{fig:entMOND}
\end{figure}


%%%%%%%%%%%%%%%%%%%%%%%%%%%%
In order to experimentally estimate the amount of entanglement one needs to measure elements of covariance matrix, which in principle require finding position and momentum correlations~\cite{Aspelmeyer-CavOpt-Rev}. The question emerges whether a simpler method could exist to detect the stronger-than-Newtonian force, e.g., by measurements of position only. This is of course possible, but it turns out that the average displacements of the masses are about three orders of magnitude smaller than their standard deviations, for the parameters considered here.
While the alternative models give rise to larger than Newtonian variances, from a practical perspective one needs make sure that the predicted increment is not mimicked by the noise in the system.
The added value of entanglement witnessing is its independence from noise characterisation, whereas the added complexities are not high since the measurements of different mechanical quadratures can be accomplished, e.g., by a choice of local phase in homodyne detection~\cite{Aspelmeyer-CavOpt-Rev}.
%The experimenter concludes the onset of modified gravity as soon as entanglement gain is certified with statistical significance.
%Finally, we would like to address the salient difference between the gravitational entanglement on considering the MOND and the DM alternatives. This work shows that the bipartite entanglement speeds up when the gravitational field is switched from Newtonian to MONDian. On the other hand, if the gravity is considered to be Newtonian, the existence of DM particles would lead to a decoherence similar to what is triggered by air molecules or solar neutrinos~\cite{QuantGrav-Decoherence-Rijavec2021,book-Decoherence-Maximilian}. While a quantitative estimation of this decoherence depends on the properties of DM particles, in any case the entanglement accumulation is slowed down, which is opposite to the impact of switching to MONDian gravity.
Our calculations assume the two particles evolve in free space, but it would be interesting to see the result of this experiment performed on Earth in free fall.
If the modification is observed in agreement with our predictions, it will imply that the Earth's gravity does not affect the bipartite system, i.e., the theories of modified gravity (in the orange region) respect SEP in locally inertial frames.
Note that SEP is certainly violated in the MOND model (black line), which motivates our next proposal.





%%%%%%%%%%%%%%%%%%%%%%%%%%%%%%%%%%%%
%%%%%%%%%%%%%%%%%%%%%%%%%%%%%%%%%%%%
\emph{Internal dynamics in an external field.}
The fundamental equations of motion in most versions of MOND are nonlinear in the gravitational field(s), implying that the internal dynamics of a bipartite system cannot be decoupled from the large external gravitational field it is embedded in (EFE)~\cite{MOND-Milogram1983,Lag4MOND-Bekenstein1984,EquiPri-Norton1985}.
It has been shown to improve galactic rotation fits~\cite{MOND-Hees2015}, and is considered an alternative to the Planet Nine hypothesis in explaining the distribution of extreme trans-Neptunian objects~\cite{MOND-P9-Migaszewski2023,MOND-P9-Mathur2023}.
The EFE asserts that MONDian modifications kick in only when the sum of \emph{all} accelerations falls below the critical level $a_0$.
In a terrestrial experiment, the bipartite system is embedded in the large gravitational field of the Earth where $g \sim 10^{11}a_0$.
The resultant bipartite dynamics is effectively Newtonian, albeit with an anisotropic Newton's constant.
In the single-field AQUAL version of MOND this is given by~\cite{EFE-Blanchet2011}:
\begin{equation}
\mathcal{G}_\theta = \frac{G}{\tilde \mu_e\sqrt{1+ 
\lambda_e \sin^2\theta}} , \hspace{1cm} : \qty{ \lambda_e =  \frac{g \tilde \mu_e'}{a_0 \tilde \mu_e} },
\label{eq:Anisotropic-G}
\end{equation}
where $\tilde \mu_e = \tilde \mu(g/a_0)$, 
$\tilde \mu_e' = \dv{\tilde \mu}{x}|_{x=g/a_0}$, and $\theta$ is the angle between the external field and the symmetry axis joining the two particles.
This motivates our proposal to verify in the laboratory the extent to which the internal attraction between the two masses depends on their orientation to the external field.
This anisotropy is tiny and undetectable with the entanglement sensitivity of current technology, hence we focus on larger masses and study their classical dynamics.


%%%%%%%%%%%%%%%%%%%%%%%%%%%%%%%%%%%%

Consider two extreme situations where the symmetry axis of the masses is either parallel or orthogonal to the external field.
Since Earth's gravity is not uniform one has to take into account the tidal effects in both orientations.
In the parallel setting the particles are located at different heights and experience a different value of $g$: the external gravity due to the Earth being weaker at the top.
%Consequently, the particle at the bottom is falling faster than the one at the top. 
Over time this tidal effect leads to an increment in the separation between the two masses.
Given that $L \ll R_E$, where $R_E$ is the radius of the Earth, the particles drift away, each having an acceleration of $\approx gL/R_E \approx 1.3 \times 10^{-11}$ m/s$^2$.
%Hence, in the considered configuration, the particles are effectively receding from each other.
In the orthogonal setting, one has to take into account the angle subtended by the symmetry axis at the center of the Earth, even for objects placed close to each other. 
This tidal force pushes the two masses closer, with each mass having an acceleration of $\approx gL/2R_E$.
Note that the internal acceleration of each particle is $G m / L^2 \approx 5.7 \times 10^{-12}$ m/s$^2$, and hence the particles are effectively moving away from (towards) each other in the parallel (orthogonal) setting. 
According to Newtonian mechanics, the ratio of time taken to undergo a relative displacement $d$ 
in parallel and orthogonal configurations is given by
\begin{equation}
    \mathcal{R}_N = \sqrt{ \frac{ gL/2R_E + G m/L^2 } { gL/R_E - G mL^2  } },
    \label{eq:ratioN}
\end{equation}
where we assume that the particles accelerate uniformly, i.e., $d \ll L$.
In accordance with Eq.~\eqref{eq:Anisotropic-G}, the same ratio in the MOND theory is
\begin{equation}
    \mathcal{R}_M = \sqrt{ \frac{ gL/2R_E + G m/L^2\tilde\mu_e\sqrt{1+\lambda_e} } { gL/R_E - G m/L^2\tilde\mu_e } }.
\end{equation}
Note, however, the following quite a subtle point. Assuming that the MOND model is the correct one, the numerical value of the gravitational constant, as established in Cavendish-like experiments, is measured using orthogonal orientation in the external field of the Earth. Therefore, we set 
$G/\tilde \mu_e \sqrt{1+\lambda_e} = \mathcal{G}_{\pi/2} = 6.67408 \times 10^{-11}$ Nm$^2$/kg$^2$, and rewrite the last ration in the form that can be directly compared with (\ref{eq:ratioN}):
\begin{equation}
    \mathcal{R}_M = \sqrt{ \frac{ gL/2R_E + \mathcal{G}_{\pi/2} m/L^2 } { gL/R_E - \sqrt{1+\lambda_e} \mathcal{G}_{\pi/2} m/L^2 } }.
    \label{eq:ratioM}
\end{equation}
It is now evident that $\mathcal{R}_M > \mathcal{R}_N$ and the ratio depends on $\lambda_e$. We therefore treat $\lambda_e$ as a free parameter and show in Fig.~\ref{fig:ratio} the difference 
$\Delta = \mathcal{R}_M - \mathcal{R}_N$.
The temporal resolution required to observe this difference in experiments is $\sim 0.1\Delta$.

\begin{figure}[!t]
    \centering
\includegraphics[width=\linewidth]{ratio.pdf}
\caption{External field effect for two masses oriented parallel and orthogonal to Earth's gravity. The vertical axis $\Delta$ shows the difference between the ratios of times to travel small distance in parallel and orthogonal configurations accroding to MOND model with parameter $\lambda_e$ (horizontal axis) and Newton's theory (tidal effects included).
The marked ``Simple'' and ``Standard'' points correspond to the two forms of interpolating function extensively used in the literature.}
    \label{fig:ratio}
\end{figure}











%%%%%%%%%%%%%%%%%%%%%%%%%%%%%%%%%%
%%%%%%%%%%%%%%%%%%%%%%%%%%%%%%%%%%
\emph{Summary.}
%%%%%%%%%%%%%%%%%%%%%%%%%%%%%%%%%%
We presented experiments with potential outcomes revealing non-Newtonian gravity. They were motivated by modifed gravity models,  and many concrete calculations were performed using MOND parameters. 
We want to emphasize that these experiments are interesting on their own, and it would be important to know their results even when they are conducted with a precision smaller than required to test MOND. 
In particular, we proposed to look at bipartite quantum and classical dynamics of microspheres. 
In the first experiment the figure of merit is their entanglement, and in the second experiment it is their gravitational attraction when they are aligned parallel or orthogonal to the external field of the Earth.





%%%%%%%%%%%%%%%%%%%%%%%%%%%%%%%%
%%%%%%%%%%%%%%%%%%%%%%%%%%%%%%%%
\begin{acknowledgements}
%%%%%%%%%%%%%%%%%%%%%%%%%%%%%%%%
This work is jointly supported by
(i) XMUM, Malaysia,
via projects \sloppy{XMUMRF/2022-C10/IPHY/0002} and \sloppy{XMUMRF/2021-C8/IPHY/0001}, 
(ii) SERB-DST, Govt.~of India, via project \sloppy{CRG/2022/009359},
(iii) project 2021/43/P/ST2/02141 co-funded by the NCN and the EU Framework Programme for Research and Innovation (Horizon 2020) under the Marie Sk\l{}odowska-Curie grant agreement 945339.
We acknowledge the National Supercomputing Mission (NSM) for providing computing resources of `PARAM Ganga' at IIT Roorkee, which is implemented by C-DAC and supported by MeitY and DST, Govt.~of India.
\end{acknowledgements}



%%%%%%%%%%%%%%%%%%%%%%%%%%%%%%%%%%%%
%%%%%%%%%%%%%%%%%%%%%%%%%%%%%%%%%%%
\appendix
\section{APPENDIX}



%%%%%%%%%%%%%%%%%%%%%%%%%%%%%
\section{Different MOND formulations}
%%%%%%%%%%%%%%%%%%%%%%%%%%%%%

The generalizations of Newton's second law of motion and Poisson's equation describe a modification to the acceleration $\vec{a}$ of an object of mass $m$ under the influence of a non-gravitational force $\vec{F}$ and gravitational force $-m\vec{\nabla}\Phi$ where $\Phi$ is determined via the solution of a set of potentially modified Poisson equations for the set of potentials $\qty{\phi_{i}}$ [see main text for details]. 
Examples of several specific proposals are given in Table~\ref{tab:ManyMOND}, where non-zero $\mu_{ij}$ and $G_{i}$ are detailed.
We note that there can additionally exist formulations of MOND where the quantity $a_{0}$ itself may depend on $\Phi$ or $|\vec{\nabla} \Phi|$~\cite{MOND2Body-Zhao2010}, where $\Phi$ appears as a `mass term' in the field equations~\cite{RMOND-Zlosnik2021}, or where terms involving higher derivatives of potentials appear in the field equations~\cite{MOND-General-Milgrom2023}.


\begin{table}
\centering
\caption{Different formulations of the MOND theory.}
\begin{tabular}{ l|c|c }
\toprule
  Model & Potentials & Properties \\
\midrule\midrule
  Modified Inertia~\cite{MOND-Milgrom1992}	& $\phi_1 = \Phi$ & 
	$\nu=\nu(|a|/a_{0})$ 	\\
&	&	$\mu_{11}=1$
\\	\midrule
  AQUAL~\cite{Lag4MOND-Bekenstein1984} & $\phi_1 = \Phi$	& 
$\nu=1$	\\
&	&	$\mu_{11}=\mu_{11}(|\vec{\nabla}\Phi|/a_{0})$ \\
& & $G_{1}=G$
\\	\midrule
  TeVeS~\cite{TeVeS-Bekenstein2004} & $(\phi_1, \phi_2) = $ &	
	$\nu=1$				\\
& $(\Phi - \varphi,\varphi)$	&	$\mu_{11} = 1$	\\
&	&	$\mu_{22} =\mu_{22}(|\vec{\nabla}\varphi|/a_{0})$
\\	
& & $G_{1}=G_{2}=G$\\
\midrule
  QUMOND ~\cite{MOND-QLinear-Milgrom2009} & $(\phi_1,\phi_2) = $ &
	$\nu=1$				\\
& $(\Phi,\varphi)$	&	$\mu_{11} = 1$	\\
&	&	$\mu_{12} =\mu_{12}(|\vec{\nabla}\varphi|/a_{0})$	\\
&	&	$\mu_{22} = 1$	\\
&	&	$G_{2}=G$
\\	\midrule
  TRIMOND ~\cite{Milgrom:2023idw} & $(\phi_1,\phi_2,\phi_3) =$ &
	$\nu=1$				\\
& $(\Phi,\varphi,\psi)$	&	$\mu_{13} = 1$	\\
&	&	$\mu_{22} =\mu_{22}(\vec{\nabla}\varphi,\vec{\nabla}\psi)$	\\
&	&	$\mu_{23} =\mu_{23}(\vec{\nabla}\varphi,\vec{\nabla}\psi)$	\\
&	&	$\mu_{31} = 1$	\\
&	&	$\mu_{32} =\mu_{32}(\vec{\nabla}\varphi,\vec{\nabla}\psi)$	\\
&	&	$\mu_{33} =\mu_{33}(\vec{\nabla}\varphi,\vec{\nabla}\psi)$	\\
&	&	$G_{1}=G$ 		
\\  \bottomrule
\end{tabular}
\label{tab:ManyMOND}
\end{table}



%%%%%%%%%%%%%%%%%%%%%%%%%%%%%
\section{Bipartite covariance matrix}
%%%%%%%%%%%%%%%%%%%%%%%%%%%%%


The covariance matrix formalism is based on the first two statistical moments of a quantum state. Given a bipartite system $AB$ with the displacement and momentum quadratures $ \hat u = ( \hat x_A, \hat p_A, \hat x_B, \hat p_B)^T$, the covariance matrix is defined as~\cite{Negativity-Vidal2002,Negativity-CV-Serafini2004,Negativity-Gaussian-Adesso2005}:
\begin{equation}
	\bm{\sigma}_{jk}  = \frac{1}{2} \ev{ \acomm{\hat u_j}{\hat u_k} }  -  \ev{ \hat u_j }\ev{ \hat u_k } 
 : \bm{\sigma} \equiv \mqty(
	\bm{\alpha} & \bm{\gamma} \\
	\bm{\gamma}^T & \bm{\beta}
	).
\end{equation}
where $\bm{\alpha}$ ($\bm{\beta}$) contains the local mode correlation for $A$ ($B$), and $\bm{\gamma}$ describes the intermodal correlation. 
In our setting the two particles have the same mass and are prepared in identical initial states, which implies that the local modes are identical at all times: $\bm{\alpha} = \bm{\beta}$.
Accordingly, $\bm{\sigma}_{22} = \bm{\sigma}_{00}$, $\bm{\sigma}_{33} = \bm{\sigma}_{11}$, $\bm{\sigma}_{23} = \bm{\sigma}_{01}$, $\bm{\sigma}_{12} = \bm{\sigma}_{03}$, 
and the rest of the elements are constrained by the symmetry property $\bm{\sigma}_{jk} = \bm{\sigma}_{kj}$.



Assuming that the force is truncated at the linear term,
the solution for the independent elements of the covariance matrix (starting from the ground state of harmonic traps) is given by~\cite{QuantGrav-Ankit2022}:
\begin{eqnarray}
	\bm{\sigma}_{00}  &=& \frac{\hbar}{4m\omega_0} \qty[ 2+\omega_0^2t^2+\qty( 1+\frac{\omega_0^2}{\omega^2} )\sinh^2(\omega t) ],   
%%%%%%%%%%%%%%%%%%%%%%%%%%%%%%%%%%%%%%%%%%%%%%
\nonumber \\
%%%%%%%%%%%%%%%%%%%%%%%%%%%%%%%%%%%%%%%%%%%%%%
	\bm{\sigma}_{02}  &=& \frac{\hbar}{4m\omega_0} \qty[ \omega_0^2t^2-\qty(1+\frac{\omega_0^2}{\omega^2})\sinh^2(\omega t) ],
%%%%%%%%%%%%%%%%%%%%%%%%%%%%%%%%%%%%%%%%%%%%%%
\nonumber \\
%%%%%%%%%%%%%%%%%%%%%%%%%%%%%%%%%%%%%%%%%%%%%%
	\bm{\sigma}_{11}   &=&  \frac{m\hbar\omega_0}{4} \qty[ 2+\qty( 1+\frac{\omega^2}{\omega_0^2} )\sinh^2(\omega t) ], 
%%%%%%%%%%%%%%%%%%%%%%%%%%%%%%%%%%%%%%%%%%%%%%
\nonumber \\
%%%%%%%%%%%%%%%%%%%%%%%%%%%%%%%%%%%%%%%%%%%%%%
	\bm{\sigma}_{13} &=& - \frac{m\hbar\omega_0}{4} \qty( 1+\frac{\omega^2}{\omega_0^2} ) \sinh^2(\omega t), 
%%%%%%%%%%%%%%%%%%%%%%%%%%%%%%%%%%%%%%%%%%%%%%
\nonumber \\
%%%%%%%%%%%%%%%%%%%%%%%%%%%%%%%%%%%%%%%%%%%%%%
	\bm{\sigma}_{01} &=& \frac{\hbar}{8} \qty[  2\omega_0 t + \qty( \frac{\omega_0}{\omega}+\frac{\omega}{\omega_0} )\sinh(2\omega t) ],   
%%%%%%%%%%%%%%%%%%%%%%%%%%%%%%%%%%%%%%%%%%%%%%
\nonumber \\
%%%%%%%%%%%%%%%%%%%%%%%%%%%%%%%%%%%%%%%%%%%%%%
	\bm{\sigma}_{03} &=& \frac{\hbar}{8} \qty[  2\omega_0 t - \qty( \frac{\omega_0}{\omega}+\frac{\omega}{\omega_0} )\sinh(2\omega t) ],
\end{eqnarray}
where $m$ is the mass, $\omega_0$ is the frequency of the trap used to prepare the initial state, and $\omega$ encodes the interaction between the two particles as follows.
For an arbitrary central force expanded in a binomial series in terms of the displacement-to-separation ratio, the parameter $\omega$ is found by equating the coefficient of the linear term with $m\omega^2/2$.
If, instead of the force, one starts with the expansion of the potential, the coefficient of the quadratic term is to be equated with $-m\omega^2/4$.
For more than one central interaction present simultaneously, the ``total'' $\omega$ is given by a Pythagoras-like sum for the individual interactions: $\omega^2 =  \omega_A^2 + \omega_B^2 + \omega_C^2 +\dots$~\cite{QuantGrav-Ankit2022}.


%%%%%%%%%%%%%%%%%%%%%%%%%%
\section{Logarithmic negativity}
%%%%%%%%%%%%%%%%%%%%%%%%%%%%%

Negativity of the partially transposed density matrix is a necessary and sufficient condition for entanglement in two--mode Gaussian states~\cite{PPTCrit-CV-Simon2000}. As a result of partial transposition, the covariance matrix is transformed to $\tilde{\bm{\sigma}}$, which differs from $\bm{\sigma}$ by a sign-flip of $\text{Det} (\bm{\gamma})$~\cite{Negativity-Gaussian-Adesso2005}.
The symplectic eigenvalues of the covariance matrix, $\tilde{\nu}_{\pm}(\bm{\sigma})$, are given by~\cite{Negativity-CV-Serafini2004,Negativity-Vidal2002}:
\begin{equation}
	\tilde{\nu}_{\pm}(\bm{\sigma}) =  \frac{1}{\sqrt{2}} \sqrt{
	 \tilde{\Sigma}(\bm{\sigma}) \pm \sqrt{\tilde{\Sigma}^2(\bm{\sigma}) - 4 \ \text{Det}\qty( \bm{\sigma} ) } }  ,
\end{equation}
where $\tilde{\Sigma}(\bm{\sigma}) = \text{Det} (\bm{\alpha}) + \text{Det} ( \bm{\beta}) - 2 \ \text{Det} ( \bm{\gamma} ) \equiv 2 \ \qty[ \text{Det} (\bm{\alpha}) - \text{Det} ( \bm{\gamma} ) ] $. Entanglement is quantified by the minimum symplectic eigenvalue via logarithmic negativity:
\begin{equation}
	E(\bm{\sigma}) = \max \Bigg[ 0, \ -\log_2\qty(  \frac{\tilde\nu_-(\bm{\sigma})}{\hbar/2} ) \Bigg].
	\label{eq:E_from_covmat}
\end{equation}





%%%%%%%%%%%%%%%%%%%%%%%%%%%%%%%%%%%%
%%%%%%%%%%%%%%%%%%%%%%%%%%%%%%%%%%%%
\bibliographystyle{apsrev4-2}
\bibliography{Bibliography}





\end{document}



